% FILE: 00_project_identity.tex
% Project: experiments-visualizer
% Thesis Role: process-intelligence-simulation-dashboard

\section{Project Identity}
\label{sec:experiments-visualizer-identity}

\subsection{Purpose}
\label{sec:experiments-visualizer-identity-purpose}

The \texttt{experiments-visualizer} is a live process intelligence simulation dashboard implemented
as a zero-dependency, pure browser-based JavaScript single-page application. Its stated purpose,
confirmed by the \texttt{README.md} and \texttt{LIVESTREAM\_STANDARDS.md} artifacts in the source
tree, is to provide a unified interactive surface where all seven subsystems of the process
intelligence doctrine stack execute simultaneously and visibly: token game simulation, A*
alignment solving, DECLARE constraint validation, EWMA concept drift detection, cryptographic
audit chaining, autonomic MAPE-K feedback control, and TypeScript boundary-law projection.

The application self-identifies its execution engine as ``WASM4PM Core JS Emulator v2.1.0'' and
is governed under the Blue River Dam Epistemic Containment Protocol, with explicit references to
doctrine files residing at \texttt{\textasciitilde/process-intelligence/doctrine/}. This
containment protocol prohibits epistemic overshoot: the dashboard may only assert what the event
evidence supports. Every rendered metric is backed by a computable formula whose inputs are
visible in the source modules.

The dashboard therefore serves two simultaneous audiences: a live demonstration surface for the
thesis claims about full-lifecycle process intelligence, and a conformance oracle that can be
inspected line-by-line to verify that no module substitutes assertion for evidence.

\subsection{Repository Surface}
\label{sec:experiments-visualizer-identity-surface}

The repository surface occupies \texttt{/Users/sac/process-intelligence/experiments/visualizer/}
and consists of 13 JavaScript modules totalling 6,947 lines, one HTML entry point
(\texttt{index.html}), one TypeScript validation file (\texttt{visualizer-validation.ts}), one
TypeScript boundary-law declaration file (\texttt{bindings.d.ts}), and supporting configuration
files (\texttt{package.json}, \texttt{tsconfig.json}). The only entry in \texttt{node\_modules/}
is the TypeScript compiler, preserving the zero-runtime-dependency contract stated in
\texttt{LIVESTREAM\_STANDARDS.md}.

The 13 JS modules decompose along subsystem boundaries:
\texttt{petrinet.js}, \texttt{token-game.js}, \texttt{alignment.js}, \texttt{astar.js},
\texttt{declare.js}, \texttt{drift.js}, \texttt{drift-detector.js}, \texttt{autonomic.js},
\texttt{blockchain.js}, \texttt{ledger.js}, \texttt{dashboard.js}, \texttt{app.js},
and the coordination layer exposed through \texttt{index.html}. Upstream checkpoint evidence
resides in \texttt{/Users/sac/process-intelligence/checkpoints/} and experiment completion
evidence in \texttt{/Users/sac/process-intelligence/experiments/checkpoint\_\_experiments\_complete.md}.

\subsection{Primary Research Role}
\label{sec:experiments-visualizer-identity-role}

Within the dissertation, the \texttt{experiments-visualizer} occupies the role of
\emph{process-intelligence-simulation-dashboard}: the artifact that makes all theoretical claims
simultaneously observable in a single browser session without external services, databases, or
network calls. It is not a toy prototype. The \texttt{LIVESTREAM\_STANDARDS.md} conformance claim
asserts v30.1.1 ultimate standard compliance, explicitly prohibiting mocks, placeholders, and
stubs. The seven subsystems are wired together through \texttt{app.js} and \texttt{dashboard.js}
so that the output of one subsystem (e.g.\ the EWMA drift score) feeds the input of another
(e.g.\ the AutonomicController conformance threshold comparison) within the same tick cycle.

This integration role makes the visualizer the primary empirical instrument for demonstrating
that prediction and coordination can co-exist in a single stateless browser runtime without
collapsing into either pure UI rendering or pure algorithmic computation.

\subsection{Key Artifacts}
\label{sec:experiments-visualizer-identity-artifacts}

The key artifacts manufactured by and for this project are:

\begin{itemize}
  \item \texttt{bindings.d.ts} — TypeScript boundary-law declarations projecting the wasm4pm
        WASM type surface (\texttt{EvidenceTs}, \texttt{EvidenceState}, \texttt{WitnessKey},
        \texttt{ReceiptShapeTs}, \texttt{GraduationCandidateTs}) into the browser runtime.
  \item \texttt{visualizer-validation.ts} — compile-time structural validation exercising all
        10 type projections with concrete values including a BLAKE3 process hash and a fitness
        score of 0.982 matching the \texttt{replay\_receipt\_sample.md} fixture.
  \item \texttt{replay\_receipt\_sample.md} — the canonical replay receipt fixture confirming
        the fitness formula $f(L,N) = 1 - \sum m / \sum c - \sum r / \sum p$ against the
        supply-chain WF-net.
  \item \texttt{checkpoint\_\_experiments\_complete.md} — confirming all 16 experiment fixtures
        verified with no deferred TODOs.
  \item \texttt{PROCESS\_INTELLIGENCE\_ALIVE\_001.md} — the GRADUATION VERIFIED verdict for the
        broader process intelligence program, covering the token game fitness formula.
\end{itemize}
